\section{Introduction}

\label{sec:introduction}

% Samsung Galaxy camera capture pipelines publish an initial image before post-processing produces the final image.
% We refer to the pipeline that produces this initial image as the draft sequence.
% For selected modes, post-processing begins only after the application enters the background, prolonging the period for which the initial image remains visible.
% We therefore added lightweight multi-frame composition to the draft sequence to reduce the visible gap to the final image.

% Each capture must produce its initial image within a fixed deadline.
% A violation of this deadline is referred to as a Capture Timeout.
% Parallel capture turns this per-capture deadline into a cross-capture control problem.
% The application may accept a new shutter request before the preceding initial image is produced, while draft sequences execute in capture order on a single worker.
% Work admitted for one capture therefore increases the waiting time of subsequent captures and reduces their remaining deadline budgets.

% The production system limits this risk by disabling optional processing in the draft sequence at a fixed overheat level.
% Thermal level, however, does not reflect queueing time, remaining work, or remaining deadline budget.
% The cutoff can suppress feasible work at or above the threshold while leaving timeout risk below it unaddressed.

% Workload adaptation and request-release control are established responses to latency and queueing pressure~\cite{approxdet-sensys,timecard-sosp,seda-sosp,breakwater-osdi}.
% In our pipeline, neither is sufficient alone: workload control cannot prevent the next capture from joining a busy worker, while release control cannot reduce work already admitted for the current capture.
% We therefore present the Budget-Aware Draft Controller, which coordinates \emph{Remaining-Sequence Admission} with \emph{Capture-Availability Pacing}.
% Admission allows an optional stage only when an upper estimate of the work remaining before initial-image publication, including mandatory encoding and saving, fits within the capture's remaining deadline budget.
% Pacing delays the next capture opportunity when projected worker occupancy would leave too little deadline budget for the next capture.
% Both decisions use observed processing durations rather than per-model timing tables or thermal-state signals.

% The two controls incur different costs.
% Pacing directly lengthens the shot-to-shot interval, whereas skipping an optional stage affects only the initial image until the final image replaces it.
% The controller therefore prioritizes responsiveness while preserving optional processing when the remaining deadline budget permits.

% Across 28 combinations of two devices, two capture conditions, and seven starting overheat levels, every run with the full controller completed the 30-capture horizon without an observed Capture Timeout.
% The unguarded baseline reached Capture Timeout in 26 of the 28 cells.
% The controller continued to execute lightweight multi-frame composition even at or above the production thermal cutoff.
% Only the full controller completed every requested capture on time in all evaluated ablation cells.

% \smallskip\noindent\textbf{Contributions.}
% We summarize our contributions below:
% \begin{itemize}[leftmargin=1.2em]
%   \item We characterize how serialized Draft execution under parallel capture creates a cross-capture deadline problem for lightweight multi-frame composition, and show why a static thermal cutoff does not reflect per-capture deadline safety.
%   \item We present the Budget-Aware Draft Controller, which coordinates remaining-sequence admission with next-capture pacing using runtime duration observations. To our knowledge, this is the first production camera controller to combine these controls under a hard per-capture deadline.
%   \item We evaluate the controller on two commercial devices across 28 device--condition--thermal combinations, examining end-to-end effectiveness, control-loop contribution, admission quality, and pacing-delay sizing.
% \end{itemize}